\documentclass[12pt,a4paper]{ctexart}
\usepackage{amsmath,amssymb,amsfonts}
\usepackage{bm}
\usepackage{graphicx}
\usepackage[margin=2.5cm]{geometry}
\usepackage{hyperref}
\usepackage{cite}
\usepackage{enumitem}
\newcommand{\R}{\mathbb{R}}
\newcommand{\abs}[1]{\lvert#1\rvert}
\newcommand{\norm}[1]{\lVert#1\rVert}

\title{基于CFD先验的AUV海流观测器}
\author{作者姓名}
\date{\today}

\begin{document}
\maketitle

\begin{abstract}
本文提出CFD先验海流观测器(PC-RCO)，结合连续Gauss-Markov正则化与自适应更新率限幅，在模型失配和传感器噪声下提供鲁棒的海流估计。验证结果表明在30\%CFD阻力扰动下PC-RCO仅退化42\%而EKF退化77\%。
\end{abstract}

\section{引言}
自主水下航行器在海洋探索中发挥着日益重要的作用。海流估计是轨迹跟踪控制中前馈补偿的关键环节。

\section{方法}
PC-RCO采用预测-校正架构，在每个时间步同时施加梯度更新和GM衰减。

\section{实验}
在XHY AUV平台上进行了系统验证，涵盖6个失配水平和4种基线方法。

\section{结论}
PC-RCO在模型失配下具有显著优于EKF的鲁棒性，连续GM正则化和自适应限幅是关键贡献。

\end{document}
